\paragraph{}
Este capítulo sintetiza as contribuições do presente trabalho, retomando os objetivos iniciais e apresentando as principais conclusões derivadas dos experimentos computacionais. Discute-se, ainda, as limitações observadas durante o desenvolvimento e as direções futuras para o aprimoramento do sistema multiagente proposto.

\section{Resumo dos resultados finais}

\paragraph{} A investigação realizada demonstrou que a arquitetura multiagente e sua variante com aprendizado por reforço apresentam desempenho consistentemente superior às meta-heurísticas clássicas aplicadas isoladamente. Ao permitir a exploração cooperativa do espaço de busca, o sistema proposto foi capaz de superar ou igualar o desempenho das linhas de base em todos os cenários de teste, confirmando a hipótese de que a inteligência coletiva atenua a sensibilidade das meta-heurísticas a configurações iniciais e topologias instáveis. No entanto, todos os métodos permanecem distantes das melhores soluções conhecidas.

\paragraph{} Os resultados agregados, corroborados pelas análises individuais por instância e, principalmente, o teste de hipóteses, evidenciam que a cooperação entre agentes especialistas permite uma otimização mais rápida e estável. Em relação aos métodos tradicionais, como \textit{SA}, \textit{GA} e \textit{TS}, a estrutura de agência proposta convergiu de maneira \textbf{mais consistente, mais rápida e mais eficaz}, validando a eficácia da abordagem híbrida no tratamento do \textit{VRPTW}.

\paragraph{} Reitera-se que, embora o sistema multiagente tenha demonstrado superioridade frente às abordagens isoladas, a implementação atual do \textit{Q-Learning} não apresentou ganhos significativos em relação à versão do MAS puro. Para além do dilema computacional clássico entre o número de \textit{steps} de aprendizado e o tempo disponível para as buscas locais dos agentes, essa equivalência decorre de limitações estruturais intrínsecas à modelagem concorrente.

\paragraph{} Primeiramente, o despacho paralelo de ações executadas pelos especialistas dentro de um mesmo ciclo de simulação faz com que a tabela de valores-\textit{Q} seja atualizada de forma diferida, ocorrendo apenas após a sincronização final do \textit{step}. Esse mecanismo quebra a semântica estritamente sequencial de transição de estados exigida pela \autoref{EqQLearning}, gerando um desalinhamento temporal entre a percepção do agente e o impacto real da ação no ambiente cooperativo.

\paragraph{} Ademais, a representação simplificada do espaço de estados por meio do par $(K, D_{\text{bin}})$ não se consolida como uma estatística suficiente do processo de otimização estocástica. Como a probabilidade de transição para uma nova solução depende diretamente da dinâmica interna e dos parâmetros probabilísticos de cada meta-heurística subjacente (como a temperatura do \textit{SA} ou a lista tabu), a hipótese de Markov mostra-se consideravelmente frágil. Sob essa perspectiva, o ambiente assume propriedades de um Processo de Decisão de Markov Parcialmente Observável (\textit{POMDP}), limitando a capacidade de convergência do algoritmo clássico sob recursos computacionais finitos.

\section{Oportunidades de melhoria}

\paragraph{} No âmbito do aprimoramento do trabalho em si, identificou-se a necessidade de evoluir a modelagem do estado no \textit{Q-Learning}. Uma proposta promissora, além de evidentemente aumentar o número de \textit{steps} e rever os parâmetros de aprendizado, é a implementação de herança de estados, permitindo que parte do valor-\textit{Q} de um estado anterior seja transmitido para o sucessor, acelerando a convergência do sistema. É possível, também, iniciar uma execução de otimização com uma tabela-\textit{Q} já existente (\textit{warm start}) de um treinamento anterior, o que aceleraria a convergência e os resultados. Adicionalmente, a otimização dos parâmetros das meta-heurísticas pode ser realizada através de abordagens teóricas ou empíricas mais refinadas.

\paragraph{} No que tange à análise macro do desempenho, o expressivo \textit{gap} absoluto de frota, discutido e quantificado na \autoref{sec:resultados}, consolida empiricamente as limitações associadas à heurística de fatiamento local. Conforme evidenciado nos resultados, a divisão gulosa sobre a representação em \textit{giant tour} atua importante fator de superestimação do número de veículos ($K$). Dessa forma, aponta-se como um trabalho futuro fundamental a substituição do fatiamento atual pela técnica de \textit{split} ótimo desenvolvida por \cite{prins2004simple}.

\paragraph{} No que tange à expansão do modelo e adições estruturais, sugere-se a hibridização profunda através de Algoritmos Meméticos, utilizando o conhecimento adquirido sobre o desempenho individual das meta-heurísticas para compensar fraquezas específicas de cada uma, como já realizado em parte pelo \textit{MAS}. Para aumentar a validade estatística do modelo, é importante testar o sistema em mais instâncias de dados com diferentes distribuições geográficas e números de pontos, evitando o viés de sobreajuste e aumentando o número de execuções por instância. Por fim, a aplicação em casos reais \cite{NVlabsOlist}, incorporando mapas georreferenciados e tempos de deslocamento não lineares, permitiria avaliar a robustez do \textit{MAS-RL} frente às imprevisibilidades do transporte urbano, aproximando a pesquisa acadêmica das demandas práticas do setor de logística.

\paragraph{} Outro ponto para o aumento da verossimilhança é o ajuste da função de custo para considerar explicitamente o tempo de espera, já que este fator representou uma parcela significativa do tempo total da rota nos experimentos realizados, impactando diretamente na eficiência logística.

\begin{figure}[H]
    \centering
    \parbox[htb]{13.0cm}
    {
        \centering
        \includegraphics[scale=0.5]{waiting_time.png}
        \caption[\small{Tempo de viagem, espera e serviço}]{\label{waiting_time} \small{Tempos gastos com viagem, espera e serviço em uma execução simples de otimização para a instância \textit{R201}.}}
    }
\end{figure}
\vspace{-0.5cm} 

\paragraph{} Nota-se, conforme ilustrado na \autoref{waiting_time}, uma desproporção crítica entre o tempo de deslocamento e a ociosidade por espera, evidenciando ineficiências que impactam a viabilidade real das rotas. Visto que o tempo de serviço total é fixo, a elevada espera indica chegadas prematuras às janelas de tempo. Esse cenário reforça o desafio de sincronização temporal do \textit{VRPTW}, validando a necessidade de mecanismos de controle para mitigar a inatividade operacional.

